\chapter{Conclusion}
\label{chap:conclusion}

This thesis investigated runtime-level scheduling for single-node multi-GPU
CUDASTF task graphs. The motivation was that GPU placement and GPU frequency
control are coupled through data movement, slack, and critical-path behavior,
but combining them into one monolithic online decision would make the
scheduler difficult to learn, explain, and stabilize. The proposed solution is
a two-layer runtime framework: a HEFT-relative residual placement scheduler
followed by a guarded, windowed proposal--commit DVFS controller.

The placement layer keeps HEFT as the latency-oriented baseline. It evaluates
candidate GPUs using data-ready time, copy time, queueing delay, and predicted
finish time, and it accepts a non-HEFT placement only when the predicted copy
saving is large enough and the normalized finish-time penalty is bounded. A
window-level contextual bandit adapts the aggressiveness of this residual gate
without directly choosing GPUs. The resulting policy remains explainable:
every accepted deviation is expressed relative to the HEFT baseline and to a
copy-saving condition.

The DVFS layer is deliberately separated from placement. Each task produces a
frequency intent, but actual GPU frequency changes are committed only at
per-device window boundaries. Criticality, backlog, heavy-task detection,
slack budget, consecutive-window requirements, cooldowns, and slowdown
guardrails prevent unstable per-task frequency switching and force recovery to
high frequency when performance risk appears.

The experimental results support the value of this layered design. With DVFS
disabled, the bandit scheduler improves geometric-mean latency by \(9.989\%\)
across the five reported benchmarks. The behavior is visible in the
locality-attribution experiment: the largest improvements coincide with large
reductions in copy bytes, while the near-neutral workload keeps the same copy
footprint and remains close to HEFT.

The final DVFS evaluation shows a smaller but meaningful energy result.
WinDVFS reduces geometric-mean energy by \(5.318\%\) with a geometric-mean
latency overhead of \(0.479\%\) under the fixed Bandit placement layer. The
geometric-mean EDP ratio improves to 0.951354, and all five reported
benchmarks remain below an EDP ratio of 1.0. The observed energy reductions
are consistent with lower average power in the selected DVFS configurations,
but the result should be interpreted as an energy--performance measurement
rather than as a full clock-residency or full-system power study.

The main conclusion is therefore that CUDASTF task-graph information can be
used to make online scheduling both more locality aware and more energy aware,
provided that the scheduler remains anchored to a strong performance baseline.
The placement contribution is the clearer and stronger result: bounded
HEFT-relative deviations can reduce unnecessary data movement and improve
runtime on copy-sensitive workloads. The DVFS contribution is more modest but
still useful: guarded window-level frequency control can reduce energy over
the benchmark execution window while keeping aggregate latency overhead small.
The end-to-end result combines these effects: Bandit+WinDVFS reduces
geometric-mean latency by \(8.795\%\), energy by \(11.874\%\), and EDP by
\(19.625\%\) relative to HEFT.

The work also identifies the limits of the current implementation. The results
are single-node measurements, the benchmark suite is limited, the criticality
and copy-cost models are lightweight approximations, and the DVFS evaluation
focuses on benchmark-window energy rather than full-system energy. The
static-replay overhead diagnostic further shows that online Bandit scheduling
cost is small at benchmark scale, but it is still a diagnostic approximation
rather than a full production-overhead study. These limitations do not
invalidate the main findings, but they define the next steps required before
the scheduler can be treated as a general-purpose production policy. Overall,
the thesis demonstrates that a carefully layered runtime scheduler can improve
locality and expose energy-saving opportunities without abandoning the
stability and interpretability of HEFT-based scheduling.
